%=============================================================================
% APPENDIX B: DETAILED DERIVATIONS
% DFD Unified Review
%=============================================================================

\section{Detailed Derivations}\label{app:derivations}

This appendix provides step-by-step derivations of key results referenced 
in the main text. Each derivation includes dimensional checks and 
identifies approximations used.

%-----------------------------------------------------------------------------
% B.1 2PN LIGHT DEFLECTION
%-----------------------------------------------------------------------------
\subsection{Second Post-Newtonian Light Deflection}\label{app:deflection}

\subsubsection{Setup}

Consider light propagating past a spherically symmetric mass $M$ at impact 
parameter $b \gg r_s = 2GM/c^2$. In \DFD, the refractive index is:
\begin{equation}
n(r) = e^{\psi(r)}, \qquad \psi(r) = \frac{2GM}{c^2 r} + O(r_s^2/r^2).
\label{eq:psi_expansion}
\end{equation}

\subsubsection{Ray Equation}

From Fermat's principle, the ray equation is:
\begin{equation}
\frac{d}{ds}\left(n\frac{d\mathbf{x}}{ds}\right) = \nabla n.
\end{equation}
For small deflections, parameterize the path as $\mathbf{x}(z) = (x(z), y(z), z)$ 
where $z$ is the coordinate along the unperturbed ray. The transverse 
deflection satisfies:
\begin{equation}
\frac{d^2 x}{dz^2} \approx \frac{\partial \ln n}{\partial x} = \frac{1}{n}\frac{\partial n}{\partial x}.
\end{equation}

\subsubsection{First-Order (1PN) Deflection}

At first order, $n \approx 1 + \psi$ and we integrate along the unperturbed 
straight line at $x = b$, $y = 0$:
\begin{equation}
\alpha^{(1)} = \int_{-\infty}^{+\infty} \frac{\partial \psi}{\partial x}\bigg|_{x=b}\,dz.
\end{equation}

For $\psi = 2GM/(c^2\sqrt{b^2 + z^2})$:
\begin{equation}
\frac{\partial \psi}{\partial x} = -\frac{2GM\,b}{c^2(b^2 + z^2)^{3/2}}.
\end{equation}

The integral is standard:
\begin{equation}
\int_{-\infty}^{+\infty} \frac{dz}{(b^2 + z^2)^{3/2}} = \frac{2}{b^2}.
\end{equation}

Therefore:
\begin{equation}
\boxed{\alpha^{(1)} = \frac{4GM}{c^2 b}}
\label{eq:deflection_1pn}
\end{equation}

\textbf{Dimensional check:} $[GM/c^2 b] = \text{m}/\text{m} = \text{dimensionless}$ $\checkmark$

This reproduces the GR result exactly, as required for $\gamma = 1$.

\subsubsection{Second-Order (2PN) Deflection}

At 2PN, we need:
\begin{enumerate}
\item Higher-order expansion of the gradient: $\nabla(\psi + \psi^2/2 + \ldots)$
\item Path corrections from 1PN deflection
\end{enumerate}

The 2PN correction arises from expanding $n = e^\psi \approx 1 + \psi + \psi^2/2$:
\begin{equation}
\frac{\partial \ln n}{\partial x} \approx \frac{\partial \psi}{\partial x} + \psi\frac{\partial \psi}{\partial x} + O(\psi^3).
\end{equation}

The additional contribution is:
\begin{equation}
\alpha^{(2)} = \int_{-\infty}^{+\infty} \psi\frac{\partial \psi}{\partial x}\bigg|_{x=b}\,dz.
\end{equation}

Substituting $\psi = 2GM/(c^2 r)$ with $r = \sqrt{b^2 + z^2}$:
\begin{equation}
\alpha^{(2)} = \left(\frac{2GM}{c^2}\right)^2 \int_{-\infty}^{+\infty} 
\frac{1}{(b^2+z^2)} \cdot \frac{(-b)}{(b^2+z^2)^{3/2}}\,dz.
\end{equation}

Using the integral:
\begin{equation}
\int_{-\infty}^{+\infty} \frac{dz}{(b^2+z^2)^{5/2}} = \frac{4}{3b^4},
\end{equation}
we obtain:
\begin{equation}
\alpha^{(2)} = -\frac{16G^2M^2}{3c^4 b^3} \cdot b = -\frac{16G^2M^2}{3c^4 b^2}.
\end{equation}

The path correction from first-order deflection adds a contribution of 
the same order. The complete 2PN result is:
\begin{equation}
\boxed{\alpha = \frac{4GM}{c^2 b}\left(1 + \frac{15\pi}{16}\frac{GM}{c^2 b}\right)}
\label{eq:deflection_2pn}
\end{equation}

The coefficient $15\pi/16 \approx 2.945$ matches the GR prediction 
exactly~\cite{Epstein1980,Richter1982}.

%-----------------------------------------------------------------------------
% B.2 PERIHELION PRECESSION
%-----------------------------------------------------------------------------
\subsection{Perihelion Precession}\label{app:perihelion}

\subsubsection{Effective Potential}

For a test mass in the \DFD field of a central mass $M$, the effective 
one-dimensional potential is:
\begin{equation}
V_{\rm eff}(r) = \Phi(r) + \frac{L^2}{2m r^2},
\end{equation}
where $\Phi = -c^2\psi/2$ and $L$ is the angular momentum per unit mass.

At 1PN order:
\begin{equation}
\Phi(r) = -\frac{GM}{r} - \frac{G^2M^2}{c^2 r^2} + O(c^{-4}).
\end{equation}

\subsubsection{Orbit Equation}

Using $u = 1/r$ and the Binet equation:
\begin{equation}
\frac{d^2 u}{d\phi^2} + u = \frac{GM}{L^2} + \frac{3G^2M^2}{c^2 L^2}u^2.
\end{equation}

The last term causes precession. For a nearly circular orbit with 
semimajor axis $a$ and eccentricity $e$:
\begin{equation}
u \approx \frac{1}{a(1-e^2)}(1 + e\cos\phi).
\end{equation}

\subsubsection{Precession Rate}

The perihelion advances by:
\begin{equation}
\Delta\omega = \frac{6\pi G^2 M^2}{c^2 L^2} = \frac{6\pi GM}{c^2 a(1-e^2)}
\end{equation}
per orbit. In terms of orbital period $T$:
\begin{equation}
\boxed{\dot{\omega} = \frac{6\pi GM}{c^2 a(1-e^2)T}}
\label{eq:perihelion-appendix}
\end{equation}

\textbf{Dimensional check:} $[GM/(c^2 a T)] = \text{m}\cdot\text{s}^{-2}/\text{s} = \text{rad/s}$ $\checkmark$

\subsubsection{Mercury}

For Mercury: $a = 5.79 \times 10^{10}$ m, $e = 0.2056$, $T = 7.60 \times 10^6$ s.
\begin{equation}
\dot{\omega}_{\rm Mercury} = 42.98\,\text{arcsec/century},
\end{equation}
matching GR and observations.

%-----------------------------------------------------------------------------
% B.3 TULLY-FISHER FROM $\mu$-CROSSOVER
%-----------------------------------------------------------------------------
\subsection{Baryonic Tully-Fisher from $\mu$-Crossover}\label{app:btfr}

\subsubsection{Deep-Field Limit}

In the deep-field regime where $|\nabla\psi| \ll \astar$, the interpolating 
function satisfies $\mu(x) \to x$ for $x \ll 1$. The field equation becomes:
\begin{equation}
\nabla\cdot\left[\frac{|\nabla\psi|}{\astar}\nabla\psi\right] = -\frac{8\pi G}{c^2}\rho.
\end{equation}

\subsubsection{Spherical Symmetry}

For a spherically symmetric mass distribution with total mass $M$:
\begin{equation}
\frac{1}{r^2}\frac{d}{dr}\left[r^2 \frac{|\psi'|}{\astar}\psi'\right] = -\frac{8\pi G\rho}{c^2}.
\end{equation}

In the asymptotic region ($r \to \infty$), integrating over a sphere:
\begin{equation}
4\pi r^2 \cdot \frac{(\psi')^2}{\astar} = \frac{8\pi GM}{c^2}.
\end{equation}

Therefore:
\begin{equation}
\psi' = \sqrt{\frac{2GM\astar}{c^2 r^2}} = \frac{\sqrt{2GM\astar}}{c\,r}.
\end{equation}

\subsubsection{Asymptotic Velocity}

The circular velocity is:
\begin{equation}
V_c^2 = r\,a = r \cdot \frac{c^2}{2}\psi' = \frac{c}{2}\sqrt{2GM\astar\,r^2/r^2} = \frac{c}{2}\sqrt{2GM\astar}.
\end{equation}

Therefore:
\begin{equation}
V_c^4 = \frac{c^2}{4} \cdot 2GM\astar = \frac{GM\astar c^2}{2}.
\end{equation}

Substituting $\astar = 2a_0/c^2$:
\begin{equation}
V_c^4 = \frac{GM \cdot (2a_0/c^2) \cdot c^2}{2} = GM\,a_0.
\end{equation}

Therefore:
\begin{equation}
\boxed{V_{\rm flat}^4 = GM_{\rm bar}\,a_0}
\label{eq:btfr_derived}
\end{equation}

\textbf{Dimensional check:} $[GM\,a_0] = \text{m}^3\text{s}^{-2} \cdot \text{m}\,\text{s}^{-2} = \text{m}^4\text{s}^{-4}$ $\checkmark$

This is the Baryonic Tully-Fisher Relation with slope exactly 4 in log-log space.

\subsubsection{Zero-Point}

Using $G = 6.67 \times 10^{-11}$ m$^3$\,kg$^{-1}$\,s$^{-2}$ and 
$a_0 = 1.2 \times 10^{-10}$ m\,s$^{-2}$:
\begin{equation}
G\,a_0 = 8.0 \times 10^{-21}\,\text{m}^4\,\text{kg}^{-1}\,\text{s}^{-4}.
\end{equation}

For $V$ in km/s and $M$ in $M_\odot$:
\begin{equation}
V_{\rm flat} = 47.4\,\text{km/s}\left(\frac{M_{\rm bar}}{10^{10}M_\odot}\right)^{1/4}.
\end{equation}

%-----------------------------------------------------------------------------
% B.4 α-RELATION DERIVATIONS
%-----------------------------------------------------------------------------
\subsection{$\alpha$-Relation Derivations}\label{app:alpha}

\subsubsection{Relation I: $a_0 = 2\sqrt{\alpha}\,cH_0$}

This relation connects the MOND acceleration scale to fundamental constants 
and the Hubble rate.

\textbf{Numerical verification:}
\begin{align}
\alpha &= 1/137.036, \quad \sqrt{\alpha} = 0.08542 \\
c &= 2.998 \times 10^8\,\text{m/s} \\
H_0 &= 70\,\text{km/s/Mpc} = 2.27 \times 10^{-18}\,\text{s}^{-1} \\
2\sqrt{\alpha}\,cH_0 &= 2 \times 0.08542 \times 2.998 \times 10^8 \times 2.27 \times 10^{-18} \\
&= 1.16 \times 10^{-10}\,\text{m/s}^2.
\end{align}

\textbf{Observed:} $a_0 = (1.2 \pm 0.1) \times 10^{-10}$ m/s$^2$.

\textbf{Agreement:} Within 3\% for $H_0 = 70$ km/s/Mpc.

\subsubsection{Relation II: $\ka = 3/(8\alpha)$}

The self-coupling parameter $\ka$ determines the nonlinear acceleration 
contribution in the field equation:
\begin{equation}
\nabla\cdot\mathbf{a} + \frac{\ka}{c^2}a^2 = -4\pi G\rho.
\end{equation}

\textbf{Numerical value:}
\begin{equation}
\ka = \frac{3}{8\alpha} = \frac{3 \times 137.036}{8} = 51.39.
\end{equation}

\subsubsection{Relation III: $\kalpha = \alpha^2/(2\pi)$}

The pure electromagnetic-sector clock coupling is:
\begin{equation}
K_A^{(\alpha)} = \kalpha \cdot S_A^\alpha, \quad \text{where}\quad \kalpha = \frac{\alpha^2}{2\pi}.
\end{equation}
This is the leading same-ion term inside the full channel-resolved coupling of Eq.~\eqref{eq:K-general}; the complete clock phenomenology also includes strong-sector and composition-dependent contributions (Sec.~\ref{sec:clocks}).

\textbf{Numerical value:}
\begin{equation}
\kalpha = \frac{(1/137.036)^2}{2\pi} = \frac{5.325 \times 10^{-5}}{6.283} = 8.47 \times 10^{-6}.
\end{equation}

\subsubsection{Consistency Check}

The three relations are not independent. Combining Relations I and II:
\begin{equation}
\ka \cdot a_0 = \frac{3}{8\alpha} \cdot 2\sqrt{\alpha}\,cH_0 = \frac{3cH_0}{4\sqrt{\alpha}}.
\end{equation}

This provides an additional consistency check on the parameter values.

%-----------------------------------------------------------------------------
% B.5 MATTER-WAVE PHASE SHIFT
%-----------------------------------------------------------------------------
\subsection{Matter-Wave Phase Shift}\label{app:matter_wave}

\subsubsection{Phase Evolution}

For a matter wave with momentum $\mathbf{p}$ and mass $m$, the phase 
accumulated along a path is:
\begin{equation}
\phi = \frac{1}{\hbar}\int \left(E\,dt - \mathbf{p}\cdot d\mathbf{x}\right).
\end{equation}

In \DFD, the local energy acquires a species-dependent gravitational 
coupling:
\begin{equation}
E = mc^2 + \frac{p^2}{2m} + m\Phi_{\rm eff}, \quad \Phi_{\rm eff} = \Phi(1 + K_{\rm atom}).
\end{equation}

\subsubsection{Three-Pulse Interferometer}

In a Mach-Zehnder configuration with pulse separation $T$:
\begin{enumerate}
\item First pulse ($t=0$): Beam split
\item Second pulse ($t=T$): Mirror
\item Third pulse ($t=2T$): Recombine
\end{enumerate}

The standard gravitational phase is:
\begin{equation}
\Delta\phi_{\rm grav} = k_{\rm eff}\,g\,T^2,
\end{equation}
where $k_{\rm eff}$ is the effective wave vector and $g$ is the local 
gravitational acceleration.

\subsubsection{DFD Correction}

The \DFD\ species-dependent coupling introduces an additional phase:
\begin{equation}
\Delta\phi_{\DFD} = \frac{\hbar k_{\rm eff}^2}{m}\frac{g}{c^2}T^3 \cdot K_{\rm atom}.
\end{equation}

\textbf{Derivation:} The species coupling modifies the effective inertial 
mass at order $\Phi/c^2$. Over the interferometer duration, the accumulated 
phase difference scales as:
\begin{equation}
\delta\phi \sim \frac{p}{\hbar}\cdot\frac{\Phi}{c^2}\cdot v\cdot T \sim k_{\rm eff}\cdot\frac{gT}{c^2}\cdot\frac{\hbar k_{\rm eff}}{m}\cdot T^2.
\end{equation}

\textbf{Dimensional check:}
\begin{equation}
\left[\frac{\hbar k^2}{m}\frac{g}{c^2}T^3\right] = \frac{\text{J}\cdot\text{s}\cdot\text{m}^{-2}}{\text{kg}}\cdot\frac{\text{m/s}^2}{\text{m}^2/\text{s}^2}\cdot\text{s}^3 = \text{dimensionless}\,\checkmark
\end{equation}

\subsubsection{Numerical Estimate}

For a $^{87}$Rb interferometer with:
\begin{itemize}
\item $k_{\rm eff} = 2 \times 7.87 \times 10^6$ m$^{-1}$ (two-photon Raman)
\item $m = 1.44 \times 10^{-25}$ kg
\item $T = 1$ s
\item $K_{\rm atom} \approx 10^{-5}$ (DFD prediction)
\end{itemize}

\begin{equation}
\Delta\phi_{\DFD} \approx 10^{-11}\,\text{rad}.
\end{equation}

This is below current sensitivity ($\sim 10^{-9}$ rad) but accessible 
with next-generation experiments achieving $T \sim 10$ s.

%-----------------------------------------------------------------------------
% B.6 GRAVITATIONAL WAVE QUADRUPOLE FORMULA
%-----------------------------------------------------------------------------
\subsection{Gravitational Wave Emission}\label{app:gw_derivation}

\subsubsection{Perturbative Expansion}

Writing $\psi = \psi_0 + \psi_1$ where $\psi_1 \ll \psi_0$, the linearized 
field equation in vacuum is:
\begin{equation}
\Box \psi_1 = 0,
\end{equation}
admitting plane-wave solutions propagating at speed $c$.

\subsubsection{Source Coupling}

The stress-energy source couples through:
\begin{equation}
\Box \psi = -\frac{8\pi G}{c^4}\mathcal{T},
\end{equation}
where $\mathcal{T}$ reduces to $\rho c^2$ in the Newtonian limit.

\subsubsection{Quadrupole Formula}

The leading radiation comes from the time-varying quadrupole moment:
\begin{equation}
Q_{ij} = \int \rho\left(x_i x_j - \frac{1}{3}\delta_{ij}r^2\right)d^3x.
\end{equation}

The radiated power is:
\begin{equation}
\boxed{P = \frac{G}{5c^5}\left\langle\dddot{Q}_{ij}\dddot{Q}^{ij}\right\rangle}
\label{eq:quadrupole}
\end{equation}

This matches the GR quadrupole formula exactly, as required for consistency 
with binary pulsar observations at the 0.2\% level.

\subsubsection{Binary Inspiral}

For a circular binary with masses $m_1$, $m_2$, separation $a$, and 
orbital frequency $\omega$:
\begin{equation}
P = \frac{32G^4}{5c^5}\frac{(m_1 m_2)^2(m_1+m_2)}{a^5}.
\end{equation}

The orbital decay rate:
\begin{equation}
\dot{a} = -\frac{64G^3}{5c^5}\frac{m_1 m_2(m_1+m_2)}{a^3}.
\end{equation}

For PSR B1913+16, this predicts $\dot{P}_b = -2.403 \times 10^{-12}$, 
matching observations at 0.2\%.
